Electricity theft is a persistent source of non-technical loss in electric-power systems. Beyond lost revenue, irregular and unbilled consumption complicates demand estimation, network planning, and the allocation of field-inspection resources. Smart meters make long consumption histories available at customer level, creating an opportunity to rank customers for investigation using data-driven models rather than relying only on manual rules. Reviews of non-technical-loss detection describe a broad progression from expert rules and engineered indicators to supervised machine learning and deep temporal models \cite{Messinis2018NTLReview}. Established supervised work also shows that consumption histories and customer-profile features contain useful discrimination signal \cite{Buzau2019NTL,Zheng2018WideDeep,Hasan2019CNNLSTM}.

The label available in the State Grid Corporation of China (SGCC) data used by AMP is nevertheless coarser than the measurements. Each customer has a binary \texttt{FLAG}, whereas the consumption record spans many days. The flag identifies a theft-associated customer; it does not identify the day on which theft occurred, its duration, or the particular 30-day window in which a verified event took place. Consequently, copying the customer flag onto every window would create propagated weak labels rather than verified window-level theft annotations. Multiple-instance learning formalizes this distinction: a labelled bag can contain instances whose individual labels are unknown \cite{Carbonneau2018MILSurvey}. For AMP, the customer is the labelled bag and temporal windows are unlabeled instances.

This distinction matters because an unusual consumption segment is not automatically theft-associated. Low usage, missingness, seasonality, occupancy changes, equipment changes, and legitimate behavioural variation can all increase reconstruction error. In the frozen AMP evidence, fully connected, convolutional, and recurrent reconstruction models learned consumption regularities but yielded weak customer-level discrimination when anomaly scores were aggregated. A later supervised diagnostic showed that some discriminative signal existed, yet scoring individually propagated windows still underused the customer label. These findings do not invalidate representation learning; they show that reconstruction abnormality and customer-level theft discrimination are different objectives.

Customer histories are also heterogeneous. Reducing all windows to one mean, one maximum, or one attention-weighted vector can erase information about within-customer variability, tails, and occasional extremes. A customer may be better characterized by the distribution of behaviours across windows than by one representative window or one pooled centroid. Deep time-series research provides many ways to encode local temporal segments \cite{Fawaz2019TimeSeriesReview}, but a suitable customer formulation must preserve useful variation without asserting that any specific window is a theft event.

The practical problem is therefore customer-level ranking under weak temporal supervision and class imbalance. The desired model must use long smart-meter histories, respect the fact that only customers are labelled, retain more than a single summary of their window behaviour, and be assessed with metrics suitable for a rare positive class. Precision--recall analysis is particularly important in imbalanced settings \cite{Saito2015PR}, while imbalance remedies and evaluation choices must be distinguished from the representation itself \cite{He2009Imbalanced}. What remains missing is a governed customer-level formulation that uses window-level temporal structure without pretending that window labels exist, preserves heterogeneity across windows, and isolates the value of learned representations from the value of aggregation itself.
